% ============================================================
% Projeto de Pesquisa — SIGA-TCC (TccManager)
% Centro Universitário de Itajubá-FEPI
%
% Baseado no template institucional (.doc) fornecido pela FEPI.
% Classe: abnTeX2 (https://www.abntex.net.br/)
% ============================================================
\documentclass[
    12pt,
    a4paper,
    oneside,
    chapter=TITLE,
    section=TITLE,
    subsection=TITLE
]{abntex2}

% ---------------------------------------------------------------
% Pacotes básicos
% ---------------------------------------------------------------
\usepackage[T1]{fontenc}
\usepackage[utf8]{inputenc}
\usepackage[brazil]{babel}
\usepackage{indentfirst}
\usepackage{graphicx}
\usepackage{microtype}
\usepackage{multirow}
\usepackage[alf]{abntex2cite}

% Fonte: o template institucional usa Arial no corpo do texto.
% "helvet" é a substituta metricamente compatível (Helvetica) usada
% para reproduzir Arial em PDF gerado por LaTeX.
\usepackage[scaled=1]{helvet}
\renewcommand{\familydefault}{\sfdefault}

% ---------------------------------------------------------------
% Margens do template institucional (medidas exatamente a partir do
% arquivo .doc original: 3cm esquerda/direita, 2,5cm superior/inferior).
% abnTeX2 é baseado na classe memoir, que usa seus próprios comandos de
% layout em vez do pacote geometry.
% ---------------------------------------------------------------
\setlrmarginsandblock{3cm}{3cm}{*}
\setulmarginsandblock{2.5cm}{2.5cm}{*}
\checkandfixthelayout

% ---------------------------------------------------------------
% Este é um documento "plano" (sem \chapter, só \section): o template
% institucional numera 1, 2, 3... e não 0.1, 0.2 (que seria o padrão
% do abnTeX2 ao combinar seção com um capítulo implícito de número 0).
% ---------------------------------------------------------------
\renewcommand{\thesection}{\arabic{section}}
\renewcommand{\thesubsection}{\thesection.\arabic{subsection}}

% ---------------------------------------------------------------
% Dados do documento — EDITE AQUI
% ---------------------------------------------------------------
% Nome do curso, usado apenas na capa (não é um campo padrão do abnTeX2).
\newcommand{\curso}{Nome do Curso}

\titulo{TÍTULO DO TRABALHO}
\autor{Nome do Aluno(a)}
\local{Itajubá}
\data{2025}
\instituicao{Centro Universitário de Itajubá-FEPI}
\orientador{Prof(a). Nome do professor}
\preambulo{Projeto de Trabalho de Conclusão de Curso.}

% ---------------------------------------------------------------
% Configurações de aparência do PDF
% ---------------------------------------------------------------
\makeatletter
\hypersetup{
  pdftitle={\@title},
  pdfauthor={\@author},
  pdfsubject={Projeto de Pesquisa},
  pdfkeywords={abnt, latex, projeto de pesquisa}
}
\makeatother

\begin{document}

% ---------------------------------------------------------------
% ELEMENTOS PRÉ-TEXTUAIS
% ---------------------------------------------------------------
\pretextual
% O template institucional numera as páginas a partir da folha de
% rosto (a capa continua sem número, mas sumário e resumo já mostram
% o número no canto superior direito — diferente do padrão do
% abnTeX2, que só numera a partir da parte textual).
\pagestyle{abntchapfirst}

% Capa customizada, reproduzindo o template institucional: nome da
% instituição e do curso aparecem acima do nome do autor (o
% \imprimircapa padrão do abnTeX2 não inclui esses campos).
\begin{capa}
  \thispagestyle{empty}
  \center
  \large\imprimirinstituicao

  \large Curso de \curso
  \vspace{2cm}

  \large\imprimirautor

  \vfill
  \begin{center}
    \bfseries\LARGE\imprimirtitulo
  \end{center}
  \vfill

  \bfseries\large\imprimirlocal

  \bfseries\large\imprimirdata
  \vspace*{1cm}
\end{capa}

% No arquivo .doc original a numeração começa em "2" na folha de
% rosto (a capa é a página 1, mas sem número impresso).
\setcounter{page}{2}

\imprimirfolhaderostostar

\tableofcontents*
\clearpage

\begin{resumo}
  Escreva um resumo na seguinte estrutura: apresentação do tema e
  objetivos da pesquisa, procedimentos metodológicos e resultados
  esperados ao término do trabalho (de 100 a 150 palavras).

  \vspace{\onelineskip}
  \noindent
  \textbf{Palavras-chave}: palavra 1. palavra 2. palavra 3.
\end{resumo}
\clearpage

% ---------------------------------------------------------------
% ELEMENTOS TEXTUAIS
% ---------------------------------------------------------------
\textual
% O template institucional não mostra o título da seção no cabeçalho,
% apenas o número da página no canto superior direito.
\pagestyle{abntchapfirst}

\section{Apresentação do tema}
\label{sec:apresentacao-tema}

% Apresente qual é a temática de seu trabalho e em qual área se encontra.
% Use linguagem clara e objetiva. Lembre-se que o tema precisa ser um
% recorte, algo específico dentro de uma área de estudo.

Apresente aqui a temática do trabalho.

\section{Justificativa}
\label{sec:justificativa}

% Justifique a importância do estudo e da temática escolhida. Diga por
% que escolheu essa temática. Use sempre linguagem formal.

Justifique aqui a importância do estudo.

\section{Problema de pesquisa}
\label{sec:problema-pesquisa}

% Problematize a temática. Ao final, o TCC precisa resolver esse
% "problema" ou indicar caminhos para a resolução. O problema é a
% principal motivação para o desenvolvimento do trabalho.

Descreva aqui o problema de pesquisa.

\section{Objetivos}
\label{sec:objetivos}

\subsection{Geral}
\label{sec:objetivo-geral}

% Trace um objetivo geral para o trabalho. Inicie com verbo no infinitivo.

Descreva aqui o objetivo geral.

\subsection{Específicos}
\label{sec:objetivos-especificos}

% Estabeleça objetivos específicos (metas menores alcançadas ao longo da
% pesquisa). Ao alcançar todos eles, o objetivo geral é alcançado.
% Inicie cada um com verbo no infinitivo — todos subordinados ao
% objetivo geral do trabalho.

\begin{enumerate}
  \item Descreva aqui o primeiro objetivo específico.
  \item Descreva aqui o segundo objetivo específico.
  \item Descreva aqui o terceiro objetivo específico.
\end{enumerate}

\section{Hipóteses}
\label{sec:hipoteses}

% Escreva sobre a(s) hipótese(s) que se espera comprovar ao final do
% trabalho.

Descreva aqui a(s) hipótese(s) da pesquisa.

\section{Referencial teórico}
\label{sec:referencial-teorico}

% Apresente os principais teóricos que serão estudados, a
% contextualização da área em que a pesquisa será desenvolvida, os
% principais conceitos e os principais trabalhos que já trataram desse
% tema.
%
% Exemplo de citação ABNT via abntex2cite:
%   \citeonline{sobrenomeAno}
%   \cite{sobrenomeAno}

Descreva aqui o referencial teórico.

\section{Procedimentos metodológicos}
\label{sec:procedimentos-metodologicos}

% Discorra sobre as etapas da pesquisa: como ela será desenvolvida.

Descreva aqui os procedimentos metodológicos.

\section{Cronograma de atividades}
\label{sec:cronograma}

% Ajuste as atividades e os meses marcados com "X" conforme o
% planejamento real do projeto.

\begin{table}[htb]
  \caption{Cronograma de atividades}
  \label{tab:cronograma}
  \centering
  \resizebox{\textwidth}{!}{%
  \footnotesize
  \begin{tabular}{|p{5.5cm}|c|c|c|c|c|c|c|c|c|c|}
    \hline
    \multirow{2}{*}{\textbf{Tarefas}} & \multicolumn{10}{c|}{\textbf{2025}} \\
    \cline{2-11}
     & Fev & Mar & Abr & Mai & Jun & Jul & Ago & Set & Out & Nov \\
    \hline
    Elaboração do tema do projeto de pesquisa & X & & & & & & & & & \\
    \hline
    Leituras relativas ao tema -- referencial teórico & & X & X & X & & & & & & \\
    \hline
    Elaboração do problema de pesquisa e hipóteses & & X & X & & & & & & & \\
    \hline
    Elaboração dos objetivos (geral e específicos) & & & X & & & & & & & \\
    \hline
    Entrega do projeto de pesquisa & & & & X & X & & & & & \\
    \hline
    Correção e revisão final & & & & & X & & & & & \\
    \hline
    Elaboração da fundamentação teórica & & & & & X & X & X & X & X & \\
    \hline
    Elaboração da metodologia e resumo & & & & & & & X & X & X & \\
    \hline
    Formatação e revisão final & & & & & & & & & & X \\
    \hline
    Correção e entrega do projeto de pesquisa finalizado & & & & & & & & & & X \\
    \hline
  \end{tabular}%
  }
  \fonte{Elaborado pelo autor (\the\year).}
\end{table}


% ---------------------------------------------------------------
% ELEMENTOS PÓS-TEXTUAIS
% ---------------------------------------------------------------
\postextual

% Título da seção de referências, igual ao template institucional
% (nesta fase do trabalho as referências ainda são preliminares).
% Precisa ser feito aqui, depois do \begin{document}, porque o pacote
% babel (idioma brazil) redefine \bibname ao carregar o idioma.
\renewcommand{\bibname}{Referências Bibliográficas Iniciais}

\bibliography{referencias}

\end{document}
